
\documentclass[../../../physics-standard_questions_all.tex]{subfiles}

\begin{document}

真空中に$xy$平面を定め，$y$軸上の点$(0, \hspace{1.2mm} L)$に正の点電荷A（電荷$Q$）を固定し，
さらに$y$軸上の点$(0, \hspace{1.2mm} -L)$に負の点電荷B（電荷$-Q$）を固定した．
クーロンの法則の比例定数を$k_0$とし，クーロン力以外の力の影響は考えない．
点$(x, \hspace{1.2mm} y)$における電場の$x$成分を$E_x(x, \hspace{1.2mm} y)$，
$y$成分を$E_y(x, \hspace{1.2mm} y)$，
電位（無限遠基準）を$\phi(x, \hspace{1.2mm} y)$と表す．

\begin{enumerate}

    \item $E_x(L, \hspace{1.2mm} L)$，$E_y(L, \hspace{1.2mm} L)$を求めよ．

    \item $E_x(x, \hspace{1.2mm} 0)$，$E_y(x, \hspace{1.2mm} 0)$，
    $\phi(x, \hspace{1.2mm} 0)$を求めよ．

    \item $E_x(0, \hspace{1.2mm} y)$，$E_y(0, \hspace{1.2mm} y)$，
    $\phi(0, \hspace{1.2mm} y)$を求めよ．
    ただし，$y \neq \pm L$とする．

    \item $y$軸に沿ってなめらかに動ける正の点電荷C（質量$m$，電荷$q$）を，
    点$(0, \hspace{1.2mm} 2L)$で静かに放したところ，Cは無限遠方へ飛び去った．
    Cが無限遠に到達したときの速さ$v_\infty$を求めよ．

\end{enumerate}

\end{document}


